\documentclass[10pt,a4paper]{amsart}
\usepackage{amsmath,amssymb,amsthm,graphicx}
% Retained minimal preamble for the Stacks Project chapter more-algebra.tex (More on Algebra).
% The chapter file pinned-original-source.tex carries only \input{preamble} on its first line, so no
% preamble text lives inside the chapter itself.  The Stacks Project uses one single shared preamble for
% every chapter; the snapshot stages it identically as sources/stacks-chow/preamble.tex and
% sources/stacks-etale-cohomology/preamble.tex.  The block below is copied mechanically, line for line,
% from sources/stacks-chow/preamble.tex lines 14-18 (the Xy-pic setup) and 154-217 (the native theorem
% declarations together with the mathematical macros).  Nothing was added, removed or redefined; the
% wrapper class and the amsmath/amssymb/amsthm/graphicx base packages are supplied by the assembler.
% Chapter: more-algebra.tex -- More on Algebra.
% For commutative diagrams we use Xy-pic
\usepackage[all]{xy}

% We use 2cell for 2-commutative diagrams.
\xyoption{2cell}

% Theorem environments.
%
\theoremstyle{plain}
\newtheorem{theorem}[subsection]{Theorem}
\newtheorem{proposition}[subsection]{Proposition}
\newtheorem{lemma}[subsection]{Lemma}

\theoremstyle{definition}
\newtheorem{definition}[subsection]{Definition}
\newtheorem{example}[subsection]{Example}
\newtheorem{exercise}[subsection]{Exercise}
\newtheorem{situation}[subsection]{Situation}

\theoremstyle{remark}
\newtheorem{remark}[subsection]{Remark}
\newtheorem{remarks}[subsection]{Remarks}

\numberwithin{equation}{subsection}

% Macros
%
\def\lim{\mathop{\mathrm{lim}}\nolimits}
\def\colim{\mathop{\mathrm{colim}}\nolimits}
\def\Spec{\mathop{\mathrm{Spec}}}
\def\Hom{\mathop{\mathrm{Hom}}\nolimits}
\def\Ext{\mathop{\mathrm{Ext}}\nolimits}
\def\SheafHom{\mathop{\mathcal{H}\!\mathit{om}}\nolimits}
\def\SheafExt{\mathop{\mathcal{E}\!\mathit{xt}}\nolimits}
\def\Sch{\mathit{Sch}}
\def\Mor{\mathop{\mathrm{Mor}}\nolimits}
\def\Ob{\mathop{\mathrm{Ob}}\nolimits}
\def\Sh{\mathop{\mathit{Sh}}\nolimits}
\def\NL{\mathop{N\!L}\nolimits}
\def\CH{\mathop{\mathrm{CH}}\nolimits}
\def\proetale{{pro\text{-}\acute{e}tale}}
\def\etale{{\acute{e}tale}}
\def\QCoh{\mathit{QCoh}}
\def\Ker{\mathop{\mathrm{Ker}}}
\def\Im{\mathop{\mathrm{Im}}}
\def\Coker{\mathop{\mathrm{Coker}}}
\def\Coim{\mathop{\mathrm{Coim}}}

% Boxtimes
%
\DeclareMathSymbol{\boxtimes}{\mathbin}{AMSa}{"02}

%
% Macros for moduli stacks/spaces
%
\def\QCohstack{\mathcal{QC}\!\mathit{oh}}
\def\Cohstack{\mathcal{C}\!\mathit{oh}}
\def\Spacesstack{\mathcal{S}\!\mathit{paces}}
\def\Quotfunctor{\mathrm{Quot}}
\def\Hilbfunctor{\mathrm{Hilb}}
\def\Curvesstack{\mathcal{C}\!\mathit{urves}}
\def\Polarizedstack{\mathcal{P}\!\mathit{olarized}}
\def\Complexesstack{\mathcal{C}\!\mathit{omplexes}}
% \Pic is the operator that assigns to X its picard group, usage \Pic(X)
% \Picardstack_{X/B} denotes the Picard stack of X over B
% \Picardfunctor_{X/B} denotes the Picard functor of X over B
\def\Pic{\mathop{\mathrm{Pic}}\nolimits}
\def\Picardstack{\mathcal{P}\!\mathit{ic}}
\def\Picardfunctor{\mathrm{Pic}}
\def\Deformationcategory{\mathcal{D}\!\mathit{ef}}
\usepackage{cleveref}
\crefname{originallemma}{original lemma}{original lemmas}
\Crefname{originallemma}{Original lemma}{Original lemmas}
\makeatletter
\newlabel{more-algebra-lemma-join-H1-regular-sequences}{{0668}{source}{}{sourcepointer-1}{}}
\newlabel{more-algebra-lemma-join-H1-regular-sequences@cref}{{[originallemma][0][]0668}{[1][0][]source}{}{}{}}
\makeatother
\title{Juxtaposing an H1-regular sequence with an H1-regular sequence modulo the ideal it generates stays H1-regular}
\author{The Stacks Project Authors}
\date{stacks-project:more-algebra.tex:local-snapshot-20261012; GNU FDL 1.2 or later}
\begin{document}
\maketitle
\setcounter{subsection}{1}
\noindent\textit{Editorial note.} All delivered mathematical text is the author's own, copied from the pinned Stacks Project chapter file more-algebra.tex; no statement or proof step was rewritten, reordered or removed.\par\medskip
\noindent\textit{Editorial note.} The chapter file pinned-original-source.tex carries only \textbackslash{}input\{preamble\} on its first line, so no preamble text lives inside the chapter. The retained file sources/stacks-more-algebra/retained-preamble.tex is a mechanical, line-for-line copy of the project's single shared preamble as staged in the snapshot, namely sources/stacks-chow/preamble.tex lines 14-18 (Xy-pic setup) and 154-217 (theorem declarations and mathematical macros). Nothing was added, removed or redefined, no package was added, and the excerpt serves only as preamble, never as a span coordinate.\par\medskip
\noindent\textit{Editorial note.} Both counted spans start at the \textbackslash{}begin line of their environment, so the statement and the author's proof are delivered whole; surrounding prose of the section is deliberately outside the counted spans and is carried as standard prerequisites instead.\par\medskip
\noindent\textit{Editorial note.} The citation fields (source\_id, source\_version, authors, source\_url, license.name, source\_title) are identical to those of the cards already published from this chapter, so that the batch selector treats them as one source instead of deferring them as a same-source clash.\par\medskip
\noindent\textit{Allowed standard prerequisites (editorial).}\begin{itemize}
\item H\_1-regularity of a sequence as the vanishing of H\_1 of its Koszul complex. used throughout the proof, in the hypothesis on f\_1 to f\_n and on the images of g\_1 to g\_m. Supplies both the goal, vanishing of H\_1 of the juxtaposed sequence, and the second correction step at proof lines 7613-7615.
\item The first Koszul differential and its kernel in degree one. used at proof lines 7588-7597 where the commutative diagram of exterior squares and A\textasciicircum{}\{n+m\} is written down. Identifies the elements of A\textasciicircum{}\{n+m\} mapping to zero as the relations for which combinations of trivial relations have to be exhibited.
\item Lifting along the quotient A to A/I and compatibility of the Koszul differential with the quotient. used at proof lines 7599-7605 to lift an element of the exterior square over A/I to A. Allows the first correction of the chosen relation, which is what pushes the last m coordinates into I.
\end{itemize}\medskip
\ifdefined\hypertarget\hypertarget{sourcepointer-1}{}\fi
\noindent\textit{Source reference.} Original lemma 0668: focused-seeds/AGP/sources/stacks-snapshot-20261011/tags.txt, source lines 7927--7927. The lemma's own canonical tag, quoted so that readers can resolve the label inside the delivered text; it points at the counted statement itself and carries no proof content.\par\medskip
\begin{lemma}
\label{lemma-join-H1-regular-sequences}
Let $A$ be a ring. Let $I$ be an ideal generated by an
$H_1$-regular sequence $f_1, \ldots, f_n$ in $A$.
Let $g_1, \ldots, g_m \in A$ be elements whose images
$\overline{g}_1, \ldots, \overline{g}_m$ form an $H_1$-regular sequence
in $A/I$. Then $f_1, \ldots, f_n, g_1, \ldots, g_m$ is an $H_1$-regular
sequence in $A$.
\end{lemma}

\begin{proof}
We have to show that $H_1(A, f_1, \ldots, f_n, g_1, \ldots, g_m) = 0$.
To do this consider the commutative diagram
$$
\xymatrix{
\wedge^2(A^{\oplus n + m}) \ar[r] \ar[d] &
A^{\oplus n + m} \ar[r] \ar[d] &
A \ar[r] \ar[d] & 0 \\
\wedge^2(A/I^{\oplus m}) \ar[r] &
A/I^{\oplus m} \ar[r] &
A/I \ar[r] & 0
}
$$
Consider an element $(a_1, \ldots, a_{n + m}) \in A^{\oplus n + m}$
which maps to zero in $A$. Because $\overline{g}_1, \ldots, \overline{g}_m$
form an $H_1$-regular sequence in $A/I$ we see that
$(\overline{a}_{n + 1}, \ldots, \overline{a}_{n + m})$ is the image
of some element $\overline{\alpha}$ of $\wedge^2(A/I^{\oplus m})$.
We can lift $\overline{\alpha}$ to an element
$\alpha \in \wedge^2(A^{\oplus n + m})$ and subtract the image of it
in $A^{\oplus n + m}$ from our element $(a_1, \ldots, a_{n + m})$.
Thus we may assume that $a_{n + 1}, \ldots, a_{n + m} \in I$.
Since $I = (f_1, \ldots, f_n)$ we can modify our element
$(a_1, \ldots, a_{n + m})$ by linear combinations of the elements
$$
(0, \ldots, g_j, 0, \ldots, 0, f_i, 0, \ldots, 0)
$$
in the image of the top left horizontal arrow to reduce to the case
that $a_{n + 1}, \ldots, a_{n + m}$ are zero. In this case
$(a_1, \ldots, a_n, 0, \ldots, 0)$ defines an element of
$H_1(A, f_1, \ldots, f_n)$ which we assumed to be zero.
\end{proof}

\end{document}
